\MathBookSourcePage{184}
\subsection{习题}
\label{sec:ch06-exercises}
\setcounter{exerciseitem}{0}

\begin{MBExercise}
\label{exr:ch06-01}
证明引理~\ref{lem:ch06-generalized-cauchy-schwarz}. 提示: 参见定理~\ref{thm:ch06-w-cycle-smoothing-property} 的证明.
\end{MBExercise}

\begin{MBExercise}
\label{exr:ch06-02}
按 W-循环和 V-循环方法访问网格的次序画出网格调度示意图. 解释这些名称的由来.
\end{MBExercise}

\begin{MBExercise}
\label{exr:ch06-03}
证明式~\eqref{eq:ch06-relaxation-operator-properties} 中的不等式. 提示: 使用 $A_k$ 在内积~\eqref{eq:ch06-mesh-dependent-inner-product} 中自伴, 以及对所有 $v\in V_k$ 有 $(A_kv,v)_k\leq\Lambda_k(v,v)_k$ 这一事实.
\end{MBExercise}

\begin{MBExercise}
\label{exr:ch06-04}
证明二重网格算法对任意 $m$ 都是收缩映射. 提示: 使用式~\eqref{eq:ch06-projected-relaxed-error-bound}.
\end{MBExercise}

\begin{MBExercise}
\label{exr:ch06-05}
当 $m$ 充分大时, 证明 W-循环方法的第 $k$ 层迭代在 $L^2$ 范数中收敛. 提示: 使用光滑性质、逼近性质和对偶论证.
\end{MBExercise}

\begin{MBExercise}
\label{exr:ch06-06}
证明引理~\ref{lem:ch06-relaxation-energy-symmetry}. 提示: 证明 $A_k$ 在内积 $a(\cdot,\cdot)$ 中自伴.
\end{MBExercise}

\begin{MBExercise}
\label{exr:ch06-07}
证明由定义~\ref{def:ch06-v-cycle-error-operator} 定义的误差算子 $E_k$($k>1$)还可写为
\[
\MathBookFormulaID{formula-ch06-exercise-07-symmetric-v-cycle-product}
\begin{aligned}
E_k={}&[(I-P_k)+R_k^mP_k][(I-P_{k-1})+R_{k-1}^mP_{k-1}]\cdots\\
&[(I-P_2)+R_2^mP_2][I-P_1][(I-P_2)+R_2^mP_2]\cdots\\
&[(I-P_{k-1})+R_{k-1}^mP_{k-1}][(I-P_k)+R_k^mP_k].
\end{aligned}
\]
\end{MBExercise}

\begin{MBExercise}
\label{exr:ch06-08}
证明具有 $m$ 个预光滑步骤的单边 V-循环方法的误差算子 $\widetilde E_k$($k>1$)可写为
\[
\MathBookFormulaID{formula-ch06-exercise-08-one-sided-v-cycle-product}
\begin{aligned}
\widetilde E_k={}&(I-P_1)[(I-P_2)+R_2^mP_2]\cdots\\
&[(I-P_{k-1})+R_{k-1}^mP_{k-1}][(I-P_k)+R_k^mP_k],
\end{aligned}
\]
而具有 $m$ 个后光滑步骤的单边 V-循环方法的误差算子是 $\widetilde E_k^*$, 即 $\widetilde E_k$ 在 $a(\cdot,\cdot)$ 中的伴随. 推出两种方法对任意 $m$ 都收敛. 提示: 使用习题~\ref{exr:ch06-07} 和定理~\ref{thm:ch06-v-cycle-level-convergence}.
\end{MBExercise}

\begin{MBExercise}
\label{exr:ch06-09}
证明具有 $m$ 个预光滑步骤的单边 W-循环方法的误差算子 $\widetilde E_k^w$ 可写为
\[
\MathBookFormulaID{formula-ch06-exercise-09-one-sided-w-cycle-factorization}
\widetilde E_k^w=F_k\widetilde E_k,
\]
其中 $\widetilde E_k$ 如习题~\ref{exr:ch06-08} 中所定义, 且 $\|F_k\|_E\leq1$. 推出单边 W-循环方法对任意 $m$ 都收敛. 提示: 使用习题~\ref{exr:ch06-08}.
\end{MBExercise}

\MathBookSourcePage{185}
\begin{MBExercise}
\label{exr:ch06-10}
证明具有 $m$ 个预光滑步骤和 $m$ 个后光滑步骤的对称 W-循环方法的误差算子 $E_k^W$ 可写为
\[
\MathBookFormulaID{formula-ch06-exercise-10-symmetric-w-cycle-factorization}
E_k^W=\widetilde E_k^*D_k\widetilde E_k,
\]
其中 $\widetilde E_k^*$ 和 $\widetilde E_k$ 如习题~\ref{exr:ch06-08} 中所定义, 且 $\|D_k\|_E\leq1$. 推出对称 W-循环方法对任意 $m$ 都收敛. 提示: 使用习题~\ref{exr:ch06-08}.
\end{MBExercise}

\begin{MBExercise}
\label{exr:ch06-11}
解式~\eqref{eq:ch06-mesh-count-recurrences}, 并建立式~\eqref{eq:ch06-dimension-asymptotic}.
\end{MBExercise}

\begin{MBExercise}
\label{exr:ch06-12}
证明引理~\ref{lem:ch06-l2-mesh-norm-equivalence} 的证明中使用的边中点求积公式对二次多项式精确. 在三维中, 若以四个面的重心为求积点, 类似结论是否成立?
\end{MBExercise}

\begin{MBExercise}
\label{exr:ch06-13}
在 $d$ 维中证明引理~\ref{lem:ch06-l2-mesh-norm-equivalence}. 提示: 使用齐次性, 并利用 Lagrange 基函数的 Gram 矩阵正定这一事实.
\end{MBExercise}

\begin{MBExercise}
\label{exr:ch06-14}
建立第 $k$ 层对称 V-循环算法的误差算子 $E_k$($k\geq2$)的如下加性表达式:
\[
\MathBookFormulaID{formula-ch06-exercise-14-additive-error-operator}
\begin{aligned}
E_k={}&\sum_{j=2}^k
R_k^m I_{k-1}^k\cdots R_{j+1}^m I_j^{j+1}
[R_j^m(Id_j-I_{j-1}^jP_j^{j-1})R_j^m]\\
&\qquad\times P_{j+1}^jR_{j+1}^m\cdots P_k^{k-1}R_k^m,
\end{aligned}
\]
其中 $P_j^{j-1}:V_j\longrightarrow V_{j-1}$ 是 Ritz 投影 $P_{j-1}$ 在 $V_j$ 上的限制, $Id_j$ 是 $V_j$ 上的恒等算子. 使用这一表达式给出命题~\ref{prop:ch06-error-operator-positive-semidefinite} 的另一证明.
\end{MBExercise}

\begin{MBExercise}
\label{exr:ch06-15}
证明第 $k$ 层多重网格迭代是线性且相容的方案, 即
\[
\MathBookFormulaID{formula-ch06-exercise-15-linear-consistent-scheme}
MG(k,z_0,g)=B_kg+C_kz_0,
\]
其中 $B_k,C_k:V_k\longrightarrow V_k$ 是线性算子, 且 $MG(k,z_0,A_kz_0)=z_0$.
\end{MBExercise}

\begin{MBExercise}
\label{exr:ch06-16}
求习题~\ref{exr:ch06-15} 中算子 $B_k$ 的递归定义. 提示: 使用第~\ref{sec:ch06-multigrid-algorithm} 节中 $MG(k,z_0,g)$ 的递归定义以及 $B_kg=MG(k,0,g)$ 这一事实.
\end{MBExercise}

\begin{MBExercise}
\label{exr:ch06-17}
证明习题~\ref{exr:ch06-15} 中的算子 $B_k$ 和 $C_k$ 满足
\[
\MathBookFormulaID{formula-ch06-exercise-17-operator-consistency}
B_kA_k+C_k=Id_k,
\]
其中 $Id_k$ 是 $V_k$ 上的恒等算子, 并推出
\[
\MathBookFormulaID{formula-ch06-exercise-17-error-factorization}
z-MG(k,z_0,g)=(Id_k-B_kA_k)(z-z_0),
\]
其中 $A_kz=g$. 提示: 使用 $MG(k,z_0,g)$ 的相容性.
\end{MBExercise}
